% Figure: PLIC reconstruction (Chapter 4). Fractions, planes, liquid polygons
% and the normal used in each cell: generated/plic.tex (generate.jl circle
% fixture, default reconstruction_method=:lhf); labels are hand-written.
\begin{figure}
  \centering
  \begin{subfigure}[b]{0.58\textwidth}
    \centering
    \begin{tikzpicture}[x=0.66cm,y=0.66cm]
      \def\fxZliquid{}%
      \begin{scope}
        \clip (3,7) rectangle (12,12);
        % Pass 1: liquid in full cells and below the PLIC lines.
        \def\fxpcell#1#2#3#4#5{\ifnum#4=1 \fill[tr liquid] (#1,#2) rectangle ++(1,1);\fi}%
        \def\fxpliquid#1{\fill[tr liquid] #1;}%
        \def\fxpseg#1#2#3#4{}%
        % Generated by docs/techref/figures/generate.jl; do not edit.
\gdef\fxCircX{7.15}\gdef\fxCircY{6.35}\gdef\fxCircR{4.3}
\fxpcell{3}{7}{0.96}{2}{1}
\fxpliquid{(3,7.450697) -- (3,7) -- (4,7) -- (4,8) -- (3.158961,8) -- cycle}
\fxpseg{3.158961}{8}{3}{7.450697}
\fxpcell{4}{7}{1.00}{1}{0}
\fxpcell{5}{7}{1.00}{1}{0}
\fxpcell{6}{7}{1.00}{1}{0}
\fxpcell{7}{7}{1.00}{1}{0}
\fxpcell{8}{7}{1.00}{1}{0}
\fxpcell{9}{7}{1.00}{1}{0}
\fxpcell{10}{7}{1.00}{1}{0}
\fxpcell{11}{7}{0.28}{2}{1}
\fxpliquid{(11,7) -- (11.427198,7) -- (11.137811,8) -- (11,8) -- cycle}
\fxpseg{11.427198}{7}{11.137811}{8}
\fxpcell{3}{8}{0.56}{2}{1}
\fxpliquid{(3.131824,8) -- (4,8) -- (4,9) -- (3.750414,9) -- cycle}
\fxpseg{3.131824}{8}{3.750414}{9}
\fxpcell{4}{8}{1.00}{1}{0}
\fxpcell{5}{8}{1.00}{1}{0}
\fxpcell{6}{8}{1.00}{1}{0}
\fxpcell{7}{8}{1.00}{1}{0}
\fxpcell{8}{8}{1.00}{1}{0}
\fxpcell{9}{8}{1.00}{1}{0}
\fxpcell{10}{8}{0.84}{2}{1}
\fxpliquid{(10,8) -- (11,8) -- (11,8.286146) -- (10.558417,9) -- (10,9) -- cycle}
\fxpseg{11}{8.286146}{10.558417}{9}
\fxpcell{11}{8}{0.02}{2}{1}
\fxpliquid{(11,8.230926) -- (11,8) -- (11.142849,8) -- cycle}
\fxpseg{11.142849}{8}{11}{8.230926}
\fxpcell{3}{9}{0.03}{2}{3}
\fxpliquid{(3.765914,9) -- (4,9) -- (4,9.287732) -- cycle}
\fxpseg{3.765914}{9}{4}{9.287732}
\fxpcell{4}{9}{0.71}{2}{1}
\fxpliquid{(4,9.287695) -- (4,9) -- (5,9) -- (5,10) -- (4.809634,10) -- cycle}
\fxpseg{4.809634}{10}{4}{9.287695}
\fxpcell{5}{9}{1.00}{1}{0}
\fxpcell{6}{9}{1.00}{1}{0}
\fxpcell{7}{9}{1.00}{1}{0}
\fxpcell{8}{9}{1.00}{1}{0}
\fxpcell{9}{9}{0.88}{2}{1}
\fxpliquid{(9,9) -- (10,9) -- (10,9.593193) -- (9.425697,10) -- (9,10) -- cycle}
\fxpseg{10}{9.593193}{9.425697}{10}
\fxpcell{10}{9}{0.16}{2}{3}
\fxpliquid{(10,9.631343) -- (10,9) -- (10.513633,9) -- cycle}
\fxpseg{10.513633}{9}{10}{9.631343}
\fxpcell{11}{9}{0.00}{0}{0}
\fxpcell{3}{10}{0.00}{0}{0}
\fxpcell{4}{10}{0.00}{2}{3}
\fxpliquid{(4.875214,10) -- (5,10) -- (5,10.0739) -- cycle}
\fxpseg{4.875214}{10}{5}{10.0739}
\fxpcell{5}{10}{0.31}{2}{1}
\fxpliquid{(5,10.089902) -- (5,10) -- (6,10) -- (6,10.52704) -- cycle}
\fxpseg{6}{10.52704}{5}{10.089902}
\fxpcell{6}{10}{0.59}{2}{1}
\fxpliquid{(6,10.51117) -- (6,10) -- (7,10) -- (7,10.669896) -- cycle}
\fxpseg{7}{10.669896}{6}{10.51117}
\fxpcell{7}{10}{0.63}{2}{1}
\fxpliquid{(7,10.668252) -- (7,10) -- (8,10) -- (8,10.583597) -- cycle}
\fxpseg{8}{10.583597}{7}{10.668252}
\fxpcell{8}{10}{0.42}{2}{1}
\fxpliquid{(8,10.594107) -- (8,10) -- (9,10) -- (9,10.24834) -- cycle}
\fxpseg{9}{10.24834}{8}{10.594107}
\fxpcell{9}{10}{0.05}{2}{1}
\fxpliquid{(9,10.269337) -- (9,10) -- (9.380233,10) -- cycle}
\fxpseg{9.380233}{10}{9}{10.269337}
\fxpcell{10}{10}{0.00}{0}{0}
\fxpcell{11}{10}{0.00}{0}{0}
\fxpcell{3}{11}{0.00}{0}{0}
\fxpcell{4}{11}{0.00}{0}{0}
\fxpcell{5}{11}{0.00}{0}{0}
\fxpcell{6}{11}{0.00}{0}{0}
\fxpcell{7}{11}{0.00}{0}{0}
\fxpcell{8}{11}{0.00}{0}{0}
\fxpcell{9}{11}{0.00}{0}{0}
\fxpcell{10}{11}{0.00}{0}{0}
\fxpcell{11}{11}{0.00}{0}{0}
\gdef\fxZC{0.31}\gdef\fxZmx{-0.304172}\gdef\fxZmy{0.695828}\gdef\fxZa{-0.133271}
\gdef\fxZnx{-0.40054}\gdef\fxZny{0.916279}
\gdef\fxZliquid{(-0.5,-0.410098) -- (-0.5,-0.5) -- (0.5,-0.5) -- (0.5,0.02704) -- cycle}
\gdef\fxZpx{0.5}\gdef\fxZpy{0.02704}\gdef\fxZqx{-0.5}\gdef\fxZqy{-0.410098}
\gdef\fxZox{0}\gdef\fxZoy{-0.191529}
\gdef\fxZfx{0.070292}\gdef\fxZfy{-0.160802}
\gdef\fxZlx{0.118093}\gdef\fxZly{-0.319953}
\gdef\fxZi{6}\gdef\fxZj{11}

        % Pass 2: cells and PLIC segments.
        \def\fxpcell#1#2#3#4#5{\draw[tr cell] (#1,#2) rectangle ++(1,1);}%
        \def\fxpliquid#1{}%
        \def\fxpseg#1#2#3#4{\draw[tr plic] (#1,#2) -- (#3,#4);}%
        % Generated by docs/techref/figures/generate.jl; do not edit.
\gdef\fxCircX{7.15}\gdef\fxCircY{6.35}\gdef\fxCircR{4.3}
\fxpcell{3}{7}{0.96}{2}{1}
\fxpliquid{(3,7.450697) -- (3,7) -- (4,7) -- (4,8) -- (3.158961,8) -- cycle}
\fxpseg{3.158961}{8}{3}{7.450697}
\fxpcell{4}{7}{1.00}{1}{0}
\fxpcell{5}{7}{1.00}{1}{0}
\fxpcell{6}{7}{1.00}{1}{0}
\fxpcell{7}{7}{1.00}{1}{0}
\fxpcell{8}{7}{1.00}{1}{0}
\fxpcell{9}{7}{1.00}{1}{0}
\fxpcell{10}{7}{1.00}{1}{0}
\fxpcell{11}{7}{0.28}{2}{1}
\fxpliquid{(11,7) -- (11.427198,7) -- (11.137811,8) -- (11,8) -- cycle}
\fxpseg{11.427198}{7}{11.137811}{8}
\fxpcell{3}{8}{0.56}{2}{1}
\fxpliquid{(3.131824,8) -- (4,8) -- (4,9) -- (3.750414,9) -- cycle}
\fxpseg{3.131824}{8}{3.750414}{9}
\fxpcell{4}{8}{1.00}{1}{0}
\fxpcell{5}{8}{1.00}{1}{0}
\fxpcell{6}{8}{1.00}{1}{0}
\fxpcell{7}{8}{1.00}{1}{0}
\fxpcell{8}{8}{1.00}{1}{0}
\fxpcell{9}{8}{1.00}{1}{0}
\fxpcell{10}{8}{0.84}{2}{1}
\fxpliquid{(10,8) -- (11,8) -- (11,8.286146) -- (10.558417,9) -- (10,9) -- cycle}
\fxpseg{11}{8.286146}{10.558417}{9}
\fxpcell{11}{8}{0.02}{2}{1}
\fxpliquid{(11,8.230926) -- (11,8) -- (11.142849,8) -- cycle}
\fxpseg{11.142849}{8}{11}{8.230926}
\fxpcell{3}{9}{0.03}{2}{3}
\fxpliquid{(3.765914,9) -- (4,9) -- (4,9.287732) -- cycle}
\fxpseg{3.765914}{9}{4}{9.287732}
\fxpcell{4}{9}{0.71}{2}{1}
\fxpliquid{(4,9.287695) -- (4,9) -- (5,9) -- (5,10) -- (4.809634,10) -- cycle}
\fxpseg{4.809634}{10}{4}{9.287695}
\fxpcell{5}{9}{1.00}{1}{0}
\fxpcell{6}{9}{1.00}{1}{0}
\fxpcell{7}{9}{1.00}{1}{0}
\fxpcell{8}{9}{1.00}{1}{0}
\fxpcell{9}{9}{0.88}{2}{1}
\fxpliquid{(9,9) -- (10,9) -- (10,9.593193) -- (9.425697,10) -- (9,10) -- cycle}
\fxpseg{10}{9.593193}{9.425697}{10}
\fxpcell{10}{9}{0.16}{2}{3}
\fxpliquid{(10,9.631343) -- (10,9) -- (10.513633,9) -- cycle}
\fxpseg{10.513633}{9}{10}{9.631343}
\fxpcell{11}{9}{0.00}{0}{0}
\fxpcell{3}{10}{0.00}{0}{0}
\fxpcell{4}{10}{0.00}{2}{3}
\fxpliquid{(4.875214,10) -- (5,10) -- (5,10.0739) -- cycle}
\fxpseg{4.875214}{10}{5}{10.0739}
\fxpcell{5}{10}{0.31}{2}{1}
\fxpliquid{(5,10.089902) -- (5,10) -- (6,10) -- (6,10.52704) -- cycle}
\fxpseg{6}{10.52704}{5}{10.089902}
\fxpcell{6}{10}{0.59}{2}{1}
\fxpliquid{(6,10.51117) -- (6,10) -- (7,10) -- (7,10.669896) -- cycle}
\fxpseg{7}{10.669896}{6}{10.51117}
\fxpcell{7}{10}{0.63}{2}{1}
\fxpliquid{(7,10.668252) -- (7,10) -- (8,10) -- (8,10.583597) -- cycle}
\fxpseg{8}{10.583597}{7}{10.668252}
\fxpcell{8}{10}{0.42}{2}{1}
\fxpliquid{(8,10.594107) -- (8,10) -- (9,10) -- (9,10.24834) -- cycle}
\fxpseg{9}{10.24834}{8}{10.594107}
\fxpcell{9}{10}{0.05}{2}{1}
\fxpliquid{(9,10.269337) -- (9,10) -- (9.380233,10) -- cycle}
\fxpseg{9.380233}{10}{9}{10.269337}
\fxpcell{10}{10}{0.00}{0}{0}
\fxpcell{11}{10}{0.00}{0}{0}
\fxpcell{3}{11}{0.00}{0}{0}
\fxpcell{4}{11}{0.00}{0}{0}
\fxpcell{5}{11}{0.00}{0}{0}
\fxpcell{6}{11}{0.00}{0}{0}
\fxpcell{7}{11}{0.00}{0}{0}
\fxpcell{8}{11}{0.00}{0}{0}
\fxpcell{9}{11}{0.00}{0}{0}
\fxpcell{10}{11}{0.00}{0}{0}
\fxpcell{11}{11}{0.00}{0}{0}
\gdef\fxZC{0.31}\gdef\fxZmx{-0.304172}\gdef\fxZmy{0.695828}\gdef\fxZa{-0.133271}
\gdef\fxZnx{-0.40054}\gdef\fxZny{0.916279}
\gdef\fxZliquid{(-0.5,-0.410098) -- (-0.5,-0.5) -- (0.5,-0.5) -- (0.5,0.02704) -- cycle}
\gdef\fxZpx{0.5}\gdef\fxZpy{0.02704}\gdef\fxZqx{-0.5}\gdef\fxZqy{-0.410098}
\gdef\fxZox{0}\gdef\fxZoy{-0.191529}
\gdef\fxZfx{0.070292}\gdef\fxZfy{-0.160802}
\gdef\fxZlx{0.118093}\gdef\fxZly{-0.319953}
\gdef\fxZi{6}\gdef\fxZj{11}

        \draw[tr exact] (\fxCircX,\fxCircY) circle (\fxCircR);
        % Pass 3: fractions and the normal used (L: LHF, W: WY fallback).
        \def\fxpcell#1#2#3#4#5{%
          \ifnum#4=2
            \node[tr value] at (#1+0.5,#2+0.5) {#3};
            \node[tr tiny,anchor=north east,text=trgrey] at (#1+0.98,#2+0.98)
              {\ifnum#5=1 L\else\ifnum#5=3 W\fi\fi};
          \fi}%
        \def\fxpseg#1#2#3#4{}%
        % Generated by docs/techref/figures/generate.jl; do not edit.
\gdef\fxCircX{7.15}\gdef\fxCircY{6.35}\gdef\fxCircR{4.3}
\fxpcell{3}{7}{0.96}{2}{1}
\fxpliquid{(3,7.450697) -- (3,7) -- (4,7) -- (4,8) -- (3.158961,8) -- cycle}
\fxpseg{3.158961}{8}{3}{7.450697}
\fxpcell{4}{7}{1.00}{1}{0}
\fxpcell{5}{7}{1.00}{1}{0}
\fxpcell{6}{7}{1.00}{1}{0}
\fxpcell{7}{7}{1.00}{1}{0}
\fxpcell{8}{7}{1.00}{1}{0}
\fxpcell{9}{7}{1.00}{1}{0}
\fxpcell{10}{7}{1.00}{1}{0}
\fxpcell{11}{7}{0.28}{2}{1}
\fxpliquid{(11,7) -- (11.427198,7) -- (11.137811,8) -- (11,8) -- cycle}
\fxpseg{11.427198}{7}{11.137811}{8}
\fxpcell{3}{8}{0.56}{2}{1}
\fxpliquid{(3.131824,8) -- (4,8) -- (4,9) -- (3.750414,9) -- cycle}
\fxpseg{3.131824}{8}{3.750414}{9}
\fxpcell{4}{8}{1.00}{1}{0}
\fxpcell{5}{8}{1.00}{1}{0}
\fxpcell{6}{8}{1.00}{1}{0}
\fxpcell{7}{8}{1.00}{1}{0}
\fxpcell{8}{8}{1.00}{1}{0}
\fxpcell{9}{8}{1.00}{1}{0}
\fxpcell{10}{8}{0.84}{2}{1}
\fxpliquid{(10,8) -- (11,8) -- (11,8.286146) -- (10.558417,9) -- (10,9) -- cycle}
\fxpseg{11}{8.286146}{10.558417}{9}
\fxpcell{11}{8}{0.02}{2}{1}
\fxpliquid{(11,8.230926) -- (11,8) -- (11.142849,8) -- cycle}
\fxpseg{11.142849}{8}{11}{8.230926}
\fxpcell{3}{9}{0.03}{2}{3}
\fxpliquid{(3.765914,9) -- (4,9) -- (4,9.287732) -- cycle}
\fxpseg{3.765914}{9}{4}{9.287732}
\fxpcell{4}{9}{0.71}{2}{1}
\fxpliquid{(4,9.287695) -- (4,9) -- (5,9) -- (5,10) -- (4.809634,10) -- cycle}
\fxpseg{4.809634}{10}{4}{9.287695}
\fxpcell{5}{9}{1.00}{1}{0}
\fxpcell{6}{9}{1.00}{1}{0}
\fxpcell{7}{9}{1.00}{1}{0}
\fxpcell{8}{9}{1.00}{1}{0}
\fxpcell{9}{9}{0.88}{2}{1}
\fxpliquid{(9,9) -- (10,9) -- (10,9.593193) -- (9.425697,10) -- (9,10) -- cycle}
\fxpseg{10}{9.593193}{9.425697}{10}
\fxpcell{10}{9}{0.16}{2}{3}
\fxpliquid{(10,9.631343) -- (10,9) -- (10.513633,9) -- cycle}
\fxpseg{10.513633}{9}{10}{9.631343}
\fxpcell{11}{9}{0.00}{0}{0}
\fxpcell{3}{10}{0.00}{0}{0}
\fxpcell{4}{10}{0.00}{2}{3}
\fxpliquid{(4.875214,10) -- (5,10) -- (5,10.0739) -- cycle}
\fxpseg{4.875214}{10}{5}{10.0739}
\fxpcell{5}{10}{0.31}{2}{1}
\fxpliquid{(5,10.089902) -- (5,10) -- (6,10) -- (6,10.52704) -- cycle}
\fxpseg{6}{10.52704}{5}{10.089902}
\fxpcell{6}{10}{0.59}{2}{1}
\fxpliquid{(6,10.51117) -- (6,10) -- (7,10) -- (7,10.669896) -- cycle}
\fxpseg{7}{10.669896}{6}{10.51117}
\fxpcell{7}{10}{0.63}{2}{1}
\fxpliquid{(7,10.668252) -- (7,10) -- (8,10) -- (8,10.583597) -- cycle}
\fxpseg{8}{10.583597}{7}{10.668252}
\fxpcell{8}{10}{0.42}{2}{1}
\fxpliquid{(8,10.594107) -- (8,10) -- (9,10) -- (9,10.24834) -- cycle}
\fxpseg{9}{10.24834}{8}{10.594107}
\fxpcell{9}{10}{0.05}{2}{1}
\fxpliquid{(9,10.269337) -- (9,10) -- (9.380233,10) -- cycle}
\fxpseg{9.380233}{10}{9}{10.269337}
\fxpcell{10}{10}{0.00}{0}{0}
\fxpcell{11}{10}{0.00}{0}{0}
\fxpcell{3}{11}{0.00}{0}{0}
\fxpcell{4}{11}{0.00}{0}{0}
\fxpcell{5}{11}{0.00}{0}{0}
\fxpcell{6}{11}{0.00}{0}{0}
\fxpcell{7}{11}{0.00}{0}{0}
\fxpcell{8}{11}{0.00}{0}{0}
\fxpcell{9}{11}{0.00}{0}{0}
\fxpcell{10}{11}{0.00}{0}{0}
\fxpcell{11}{11}{0.00}{0}{0}
\gdef\fxZC{0.31}\gdef\fxZmx{-0.304172}\gdef\fxZmy{0.695828}\gdef\fxZa{-0.133271}
\gdef\fxZnx{-0.40054}\gdef\fxZny{0.916279}
\gdef\fxZliquid{(-0.5,-0.410098) -- (-0.5,-0.5) -- (0.5,-0.5) -- (0.5,0.02704) -- cycle}
\gdef\fxZpx{0.5}\gdef\fxZpy{0.02704}\gdef\fxZqx{-0.5}\gdef\fxZqy{-0.410098}
\gdef\fxZox{0}\gdef\fxZoy{-0.191529}
\gdef\fxZfx{0.070292}\gdef\fxZfy{-0.160802}
\gdef\fxZlx{0.118093}\gdef\fxZly{-0.319953}
\gdef\fxZi{6}\gdef\fxZj{11}

      \end{scope}
      \draw[tr block,line width=1.1pt] (\fxZi-1,\fxZj-1) rectangle ++(1,1);
      \node[tr small,anchor=south] at (\fxZi-0.5,12) {$c$};
    \end{tikzpicture}
    \caption{Exact interface (dashed), PLIC segments $\interface_c$ and liquid
      parts; numbers: $\condfrac{\liq}{c}$; L/W: LHF or WY normal.}
    \label{fig:interface:plic:patch}
  \end{subfigure}%
  \hfill
  \begin{subfigure}[b]{0.39\textwidth}
    \centering
    \begin{tikzpicture}[x=3.3cm,y=3.3cm]
      \fill[tr liquid] \fxZliquid;
      \draw[tr block] (-0.5,-0.5) rectangle (0.5,0.5);
      \draw[tr plic] (\fxZpx,\fxZpy) -- (\fxZqx,\fxZqy);
      % Perpendicular from the centroid to the line: the thesis' distance.
      \draw[black!70,dash pattern=on 1.5pt off 1.2pt,line width=0.5pt] (0,0) -- (\fxZfx,\fxZfy);
      \node[tr small,anchor=west,inner sep=2pt] at ($(0,0)!0.5!(\fxZfx,\fxZfy)$) {$\abs{\idist{c}}$};
      \node[tr centroid] at (0,0) {};
      \node[tr small,anchor=south west] at (0,0) {$\pos{c}$};
      \node[circle,fill=trpurple,inner sep=0pt,minimum size=3.6pt] (o) at (\fxZox,\fxZoy) {};
      \node[tr small,anchor=north,text=trpurple,inner sep=2.5pt] at (o.south) {$\icentroid{c}$};
      \draw[tr vector] (o) -- ++(0.3*\fxZnx,0.3*\fxZny) node[anchor=south west,tr label] {$\inormal_c$};
      \node[tr label,anchor=south east] at (\fxZqx,\fxZqy) {$\lnormal{c}\cdot\lcoord=\lconst{c}$};
      \node[tr label] at (0.12,-0.33) {$\localpart{\liq}{c}$};
      \node[tr small,anchor=north west] at (-0.5,-0.5) {$\refcell=[-\half,\half]^2$};
    \end{tikzpicture}
    \caption{Cell $c$ in local coordinates; $\condfrac{\liq}{c}=\chkPlicC$.}
    \label{fig:interface:plic:cell}
  \end{subfigure}
  \caption{PLIC reconstruction~\cref{eq:interface:plane,eq:interface:area,eq:interface:intercept}
    on the circle fixture of \code{figures/generate.jl} (unit cells, radius
    $4.3$, liquid inside; window $3\le\coordk{1}\le12$, $7\le\coordk{2}\le12$).
    Each mixed cell carries one half-plane; the segments are not connected
    across faces. The plane constant $\lconst{c}$ is fixed by the area of
    $\localpart{\liq}{c}$; with square cells $\inormal_c$ is parallel to
    $\lnormal{c}$ and $\idist{c}=\lconst{c}/\abs{\lnormal{c}}_2=\chkPlicDistance$ in
    units of the cell width. The midpoint $\icentroid{c}$ of the segment is the
    origin of the curvature fit of \cref{sec:interface:curvature}.}
  \label{fig:interface:plic}
\end{figure}
